\chapter{Benchmark Algorithms}\label{ch:mx:benchmark-algorithms}

Most institutional equity orders are handed to algorithms that promise a benchmark: the day's volume-weighted price, the arrival price, the close.
Each promise is a schedule, and each schedule is only as good as its forecast of the day's volume. This chapter builds the five classic algorithms on
chapter~14's scheduler interface, forecasts intraday volume three ways, and runs everything in the simulated market to measure how much of an
algorithm's miss against its benchmark is the volume forecast's fault, what a \omterm{def:mx:benchmark-algorithms:vwap}{participation algorithm} does to the volume it follows, and how a
benchmark can be moved by the order measured against it.

\section{The algorithms and their benchmarks}

\begin{definition}[Execution algorithm]\label{def:mx:benchmark-algorithms:algo}
An \emph{execution algorithm}\index{execution algorithm} is a program that works a \omterm{def:mx:the-almgren-chriss-framework:parent}{parent order} for a client or a desk: it decides how much to trade
in each interval (the schedule) and sends the \omterm{def:mx:the-almgren-chriss-framework:parent}{child orders}, measured against a benchmark price named in advance.
\end{definition}

\begin{definition}[VWAP, TWAP and participation algorithms]\label{def:mx:benchmark-algorithms:vwap}
A \emph{VWAP algorithm}\index{VWAP algorithm} trades each interval's quantity in proportion to the share of the day's volume it expects there, aiming
at the day's volume-weighted average price. A \emph{TWAP algorithm}\index{TWAP algorithm} trades equal quantities in equal intervals, aiming at the
time-weighted average price. A \emph{participation algorithm}\index{participation algorithm} trades a fixed percentage of volume as the volume prints,
so that its schedule follows the market instead of a forecast.
\end{definition}

\begin{definition}[Implementation-shortfall and target-close algorithms]\label{def:mx:benchmark-algorithms:is}
An \emph{implementation-shortfall algorithm}\index{implementation-shortfall algorithm} trades on an optimal-execution schedule (chapter~14's
Almgren--Chriss trajectory, front-loaded by the risk aversion), aiming at the arrival price. A \emph{target-close algorithm}\index{target-close
algorithm} aims at the closing price: it trades part of the order in the closing auction with a market-on-close order and the rest in the last
part of the day.
\end{definition}

The benchmarks are those of One Quant Book 7 (chapters 2, 19 and 23): the volume-weighted average price, the arrival price for implementation
shortfall, the closing price. Madhavan's encyclopaedia entry describes the VWAP as a common benchmark for judging trading and execution, and Kissell's
textbook (2014) catalogues the algorithms brokers built on it. Each algorithm is easy to measure against its own benchmark, and that is the trap:
the benchmark is not what the client pays against the price the order met.

In the simulated market (\texttt{firm.agentmkt} on \texttt{firm.exchsim}, a day of 26 one-minute bins whose activity follows a U-shaped curve), each
algorithm buys 15\,000 shares, about 5\% of an ordinary day's volume, on the same eight days with the same order flow, and each day is also run
without it. Prices are in ticks, one cent on a USD~100 stock, so a tick is a basis point; positive numbers are what the buyer paid above the
reference. \Cref{tab:mx:benchmark-algorithms:five} reports the miss against each algorithm's own benchmark, how far its buying moved that benchmark
(the benchmark with the algorithm minus the same day's without it), the impact cost (fills against the mid the same day had at the same instants
without the algorithm), and the cost against the arrival price.

\begin{table}[ht]
\centering
\small
\begin{tabular}{@{}lrrrr@{}}
\toprule
algorithm & vs benchmark (s.d.) & benchmark moved & impact (s.e.) & vs arrival (s.e.) \\
\midrule
VWAP & 0.6 (0.2) & 2.0 & 2.7 (0.4) & 3.0 (2.1) \\
TWAP & 0.8 (0.5) & 2.1 & 2.9 (0.2) & 3.3 (2.2) \\
participation 10\% & 1.5 (0.7) & 3.2 & 4.7 (0.6) & 5.3 (1.7) \\
shortfall, $\kappa T=4$ & 2.9 (5.1) & 0.0 & 2.9 (0.5) & 2.9 (1.8) \\
target close & $-5.9$ (3.6) & 12.6 & 6.8 (1.1) & 7.3 (2.3) \\
\bottomrule
\end{tabular}
\caption{The five algorithms buying 15,000 shares on the same eight simulated days (ticks per share); the \omterm{def:mx:benchmark-algorithms:is}{target-close algorithm} sends 30\% to
the closing auction. Data:
\texttt{mx\_algos.benchmark\_study}.}\label{tab:mx:benchmark-algorithms:five}
\end{table}

Three things show. The \omterm{def:mx:benchmark-algorithms:vwap}{VWAP algorithm} misses the VWAP by 0.6 ticks with a standard deviation of 0.2, the tightest of all; but its own buying raised
the day's VWAP by 2.0 ticks, so the cost it caused, 2.7 ticks of impact, is more than four times its reported slippage. The \omterm{def:mx:benchmark-algorithms:is}{target-close algorithm}
``beats'' the close by 5.9 ticks because its 4\,500-share market-on-close order, crossed against a thin closing book, raised the closing price by 12.6
ticks: it moved its own benchmark. Against the arrival price, the day's own price moves (a standard deviation of about six ticks) swamp the
differences on eight days; only the common-random-number impact cost ranks the algorithms reliably, and there the participation and close
algorithms, which trade fastest or in one block, cost the most.

\section{Forecasting intraday volume}

\begin{definition}[Intraday volume forecast]\label{def:mx:benchmark-algorithms:forecast}
An \emph{intraday volume forecast}\index{intraday volume forecast} gives, before the open or as the day unfolds, the expected volume of each interval
of the day, or its share of the day's total; the intraday volume profile (One Quant Book 7, chapter~5) is its simplest form.
\end{definition}

The simulated panel has twenty stocks over 250 days of 26 bins. Each stock has its own U-shaped profile with a run-up into the close; the volume of a
bin multiplies the stock's level by a market factor (persistent from day to day and autocorrelated within the day), a day-level shock, a stock-specific
AR(1) deviation within the day, and, on 8\% of stock-days, news that makes the volume surge from a random bin and fade over about four bins. These
curves are \texttt{firm.agentmkt}'s activity: they scale the rate of its \omterm{def:mx:why-there-is-a-spread:informed}{noise traders}' orders and of its liquidity providers' quotes, and the volume
a day prints is its \omterm{def:mx:why-there-is-a-spread:informed}{noise traders}' orders, which the simulator draws in advance from the day's seed. \Cref{fig:mx:benchmark-algorithms:profile} shows a
stock's average profile, an ordinary day and a news day.

\begin{omfigure}
\begin{tikzpicture}
\begin{axis}[width=11cm,height=5.4cm,xlabel={\small bin (one-minute intervals of the simulated day)},ylabel={\small share of the day's volume (\%)},
  tick label style={font=\footnotesize},grid=major,grid style={gray!25},xmin=0.5,xmax=26.5,ymin=0,ymax=16,
  legend columns=3,legend style={font=\footnotesize,at={(0.5,-0.3)},anchor=north,draw=none,fill=none,
  /tikz/every even column/.append style={column sep=8pt}},table/col sep=comma]
\addplot[omDef,very thick] table[x=bin,y=typical]{figdata/microstructure/16-benchmark-algorithms/profile.csv};
\addplot[omThm,thick,mark=*,mark size=1.2pt] table[x=bin,y=normal]{figdata/microstructure/16-benchmark-algorithms/profile.csv};
\addplot[omProp,thick,mark=square*,mark size=1.2pt] table[x=bin,y=news]{figdata/microstructure/16-benchmark-algorithms/profile.csv};
\legend{average over 250 days,an ordinary day,a news day}
\end{axis}
\end{tikzpicture}

\omcaption{Intraday volume shares of one simulated stock: its average profile, the U with a run-up to the close, and two single days. On the news day
the volume surges in bin 21, so every earlier bin's share of the day falls below the profile. Data: \texttt{mx\_algos.flow}.}\label{fig:mx:benchmark-algorithms:profile}
\end{omfigure}

Five forecasters are compared on the 190 days after a 60-day training period, for every stock:
\begin{itemize}
\item the \emph{static curve}: the average profile of the 60 training days, never updated;
\item the \emph{rolling} profile: the average of the last 20 days' shares;
\item \emph{level plus AR}: the rolling profile with a model of the day, $\log v_k=\log(\bar V p_k)+\ell+e_k$, where the level $\ell$ lasts all day and
the deviation $e_k$ is AR(1); after $k$ bins the level is estimated by generalised least squares from the deviations seen and the last deviation is
let fade (\cref{lst:mx:benchmark-algorithms:levelar});
\item \emph{PCA-ARMA}: Białkowski, Darolles and Le Fol's (2008) decomposition of turnover into a common component, extracted by principal component
analysis (One Quant Book 4, chapter~22) across the stocks and forecast by its average over the last 20 days, and a stock-specific component forecast
by an ARMA model (they also used threshold autoregressions); here AR(1);
\item the same, \emph{dynamic}: the specific component re-forecast after each bin from the last one observed.
\end{itemize}

\omcode{code/firm/algos/firm_algos.py}{65}{85}{The level-plus-AR volume model: its five parameters are fitted on a history of days, and after each bin it re-estimates the day's level by generalised least squares and lets the last deviation fade.}\label{lst:mx:benchmark-algorithms:levelar}

A forecast error matters only through what it does to the schedule's price. Write the schedule's shares $a_k$ and the day's volume shares $b_k$, and
$D_k=\sum_{i\le k}(a_i-b_i)$ their cumulative difference.

\begin{proposition}[What the volume forecast costs]\label{prop:mx:benchmark-algorithms:tracking}
If the bin VWAPs move by independent increments of variance $s^2$, independent of the schedule's errors, the part of the VWAP slippage due to the
schedule, $\sum_k(a_k-b_k)w_k$, has variance $s^2\,\E\sum_{k<B}D_k^2$.
\end{proposition}

\begin{proof}
Since $\sum_k(a_k-b_k)=0$, summation by parts gives $\sum_k(a_k-b_k)w_k=-\sum_{k<B}D_k(w_{k+1}-w_k)$, a sum of independent increments weighted by
$D_k$.
\end{proof}

The tracking error $s\sqrt{\sum D_k^2}$ is what \cref{fig:mx:benchmark-algorithms:forecast} reports, for a stock with 2\% daily volatility, on the 3\,800
stock-days of the test period and on the 331 of them (8.7\%) whose volume exceeded 1.5 times the stock's average of the previous 20 days. Konishi
(2002) solved the static version of the problem, the slicing that minimises this variance when volume and price are random.

\begin{omfigure}
\begin{tikzpicture}
\begin{axis}[width=11cm,height=5.6cm,ybar,bar width=9pt,ylabel={\small VWAP tracking error (basis points)},
  xtick={0,1,2,3,4},xticklabels={{static\\curve},{rolling},{level\\+ AR},{PCA-\\ARMA},{PCA-ARMA\\dynamic}},
  xticklabel style={align=center},
  tick label style={font=\footnotesize},ymin=0,ymax=24,grid=major,grid style={gray!25},enlarge x limits=0.14,
  legend columns=2,legend style={font=\footnotesize,at={(0.5,-0.22)},anchor=north,draw=none,fill=none,
  /tikz/every even column/.append style={column sep=8pt}},table/col sep=comma]
\addplot[fill=omDef!70,draw=omDef] table[x=i,y=normal]{figdata/microstructure/16-benchmark-algorithms/forecast.csv};
\addplot[fill=omProp!70,draw=omProp] table[x=i,y=high]{figdata/microstructure/16-benchmark-algorithms/forecast.csv};
\legend{other days,high-volume days}
\end{axis}
\end{tikzpicture}

\omcaption{The schedule part of the VWAP slippage (root mean square of $s\sqrt{\sum D_k^2}$, a stock with 2\% daily volatility) for the five
forecasters: 3,469 other stock-days and 331 high-volume ones. The level-plus-AR and PCA-ARMA dynamic schedules update during the day. Data:
\texttt{mx\_algos.forecast\_study}.}\label{fig:mx:benchmark-algorithms:forecast}
\end{omfigure}

Three lessons. Forecasting shares is hard: the best forecaster's total error in the day's shares, $\sum_k|a_k-b_k|$, is still 0.306 (against 0.324 for
the static curve), because most of a bin's deviation is noise nobody can predict. High-volume days cost more, 17.7 basis points of tracking error for
the static curve against 13.7 on other days, because their extra volume comes in surges that reshape the day. And updating helps a little on ordinary
days (the PCA-ARMA dynamic schedule gives 12.7 against 13.7) but can hurt on the days that matter: its additive AR(1), fitted mostly on quiet days,
mistakes a news surge for noise that will revert and ends at 20.6 basis points, worse than not updating. The level-plus-AR model, which separates a
day's level from its transient deviations, is the only one to do better on both kinds of day (13.3 and 17.2). The rolling profile is no better than
the static curve here because the simulated profile does not drift; on real data it does.

\section{VWAP: static and dynamic}

\begin{definition}[Dynamic VWAP]\label{def:mx:benchmark-algorithms:dvwap}
A \emph{dynamic VWAP}\index{dynamic VWAP} schedule recomputes, after each interval, the quantity left over the intervals left in proportion to an
updated forecast of their volume, instead of fixing the whole schedule before the open.
\end{definition}

The slippage of a buy against the day's VWAP splits exactly. With $q_k$ shares bought in bin $k$ at an average $f_k$, the bin's market volume $v_k$
and VWAP $w_k$, $Q=\sum q_k$ and $V=\sum v_k$,
\[
\frac{\sum_kq_kf_k}{Q}-\frac{\sum_kv_kw_k}{V}=\underbrace{\sum_k\frac{q_k}{Q}(f_k-w_k)}_{\text{execution}}+\underbrace{\sum_k\Bigl(\frac{q_k}{Q}-\frac{v_k}{V}\Bigr)w_k}_{\text{schedule}}:
\]
the execution part is what each bin's \omterm{def:mx:the-almgren-chriss-framework:parent}{child orders} paid against that bin's VWAP (the spread and the impact), the schedule part is the volume
forecast's error priced by the day's moves. \texttt{firm\_algos.decompose} computes both; the tests check that they add up to the slippage to the
cent.

\omcode{code/microstructure/16-benchmark-algorithms/python/mx_algos.py}{283}{290}{The static and dynamic VWAP schedules handed to the simulated execution agent: the dynamic one re-forecasts the bins left from the other traders' volume so far and gives the next bin its share of what is left to buy.}\label{lst:mx:benchmark-algorithms:schedules}

Humphery-Jenner (2011) built a \omterm{def:mx:benchmark-algorithms:dvwap}{dynamic VWAP} that lets the schedule react to news during the day. Here both schedules buy 15\,000 shares (4.95\% of the volume on
the ordinary days, 2.5\% on the busy ones) on sixteen high-volume and sixteen other stock-days of the simulated panel, in the simulated market, with
the day's flow the same for both.
\Cref{fig:mx:benchmark-algorithms:vwap} plots each day's slippage against its schedule part.

\begin{omfigure}
\begin{tikzpicture}
\begin{axis}[width=8.6cm,height=6.4cm,xlabel={\small schedule part (ticks per share)},ylabel={\small slippage against the VWAP (ticks)},
  tick label style={font=\footnotesize},grid=major,grid style={gray!25},xmin=-3,xmax=8,ymin=-3,ymax=8,
  legend columns=3,legend style={font=\footnotesize,at={(0.5,-0.26)},anchor=north,draw=none,fill=none,
  /tikz/every even column/.append style={column sep=6pt}},table/col sep=comma]
\addplot[omDef,only marks,mark=*,mark size=1.8pt] table[x=sched,y=total]{figdata/microstructure/16-benchmark-algorithms/vwap_high_static.csv};
\addplot[omProp,only marks,mark=square*,mark size=1.8pt] table[x=sched,y=total]{figdata/microstructure/16-benchmark-algorithms/vwap_high_dynamic.csv};
\addplot[gray,only marks,mark=o,mark size=1.5pt] table[x=sched,y=total]{figdata/microstructure/16-benchmark-algorithms/vwap_normal_static.csv};
\addplot[black,dashed,domain=-3:8] {x};
\legend{static (high volume),dynamic (high volume),static (other days)}
\end{axis}
\end{tikzpicture}

\omcaption{VWAP slippage of a 15,000-share buy against its schedule part, one point per simulated day. The dashed line is slippage equal to the
schedule part; the points sit about 0.7 tick above it, the execution part's spread and impact. Data: \texttt{mx\_algos.vwap\_study}.}\label{fig:mx:benchmark-algorithms:vwap}
\end{omfigure}

On the high-volume days the schedule part explains 99.1\% of the variance of the static VWAP's slippage and 98.6\% of the dynamic one's: the
slippage's standard deviation is 2.24 ticks for the static schedule and 1.75 for the dynamic one, the schedule part's 2.17 and 1.71, while the
execution part varies by only 0.21 and 0.20 around a mean of 0.66 and 0.73 ticks (the half-spread and impact every child pays). On the other days the
slippage is smaller (standard deviations 0.48 and 0.41) and the schedule part explains less of it, 73.9\% and 68.5\%. The dynamic schedule's lower
dispersion on high-volume days is not a reliable gain on sixteen days: it did better on seven of them, the paired difference in squared slippage
is 1.5 with a standard error of 1.8, and most of it comes from one news day (7.5 ticks for the static schedule, 5.4 for the dynamic). The panel's
3\,800 days, where the tracking error falls from 17.7 to 17.2 basis points, measure the same effect precisely: real, and small.

\section{Participation and its feedback}

A \omterm{def:mx:benchmark-algorithms:vwap}{participation algorithm} trades in each bin a share $r$ of the previous bin's volume. The volume it sees includes its own prints, and that makes it
chase itself: if it trades $q$ against the other traders' $o$ in each bin, then $q=r(o+q)$, so it trades $q/o=r/(1-r)$ of the volume the market
would have had without it, 25\% at $r=20\%$. Two such algorithms buying the same stock each count the other: $q=r(o+2q)$, so each trades
$r/(1-2r)$, 33\%, and the total volume grows by $1/(1-2r)$. And if the other traders react to its buying (here the fundamentalists sell into the
price it pushes up), their extra volume feeds the same loop.

\omcode{code/microstructure/16-benchmark-algorithms/python/mx_algos.py}{328}{335}{A participation schedule for the simulated agent: each bin, a share of the last bin's whole volume, its own prints included, until the order is done.}\label{lst:mx:benchmark-algorithms:pov}

\Cref{tab:mx:benchmark-algorithms:pov} runs one and two such buyers of 15\,000 shares at $r=20\%$, on a normal day and on a thin one with 40\% of the
normal activity, eight days each, against the same days without them.

\begin{table}[ht]
\centering
\small
\begin{tabular}{@{}llrrrrr@{}}
\toprule
day & buyers & minutes & participation & pace & others' extra volume (s.e.) & impact cost (s.e.) \\
\midrule
normal & one & 5.2 & 17.0\% & 21.1\% & 0.12 (0.10) & 6.4 (0.7) \\
normal & two & 4.6 & 15.9\% & 24.6\% & 0.12 (0.04) & 10.0 (0.9) \\
thin & one & 11.6 & 19.7\% & 24.9\% & 0.04 (0.06) & 5.4 (0.8) \\
thin & two & 9.5 & 19.2\% & 29.4\% & $-0.05$ (0.04) & 5.7 (0.6) \\
\bottomrule
\end{tabular}
\caption{\omterm{def:mx:benchmark-algorithms:vwap}{Participation algorithms} at 20\% on normal and thin simulated days: time to buy 15,000 shares, participation (own over all volume),
pace (own over the volume the same day had without the algorithms), the other traders' extra volume per share bought, and the impact cost in ticks.
For two buyers the numbers are the first one's. Data: \texttt{mx\_algos.pov\_study}.}\label{tab:mx:benchmark-algorithms:pov}
\end{table}

The participation stays near its target; the pace does not. On the normal day the order lasts five minutes, too short for the loop to build (the
last bin is also partial), and the pace is 21.1\%; on the thin day it lasts almost twelve and reaches 24.9\%, the $r/(1-r)$ of the argument. With a
second buyer the pace rises to 24.6\% and 29.4\%, toward $r/(1-2r)$, and on the normal day the impact cost rises from 6.4 to 10.0 ticks: each algorithm
pays for the other's buying and trades faster because of it. The other traders' reaction adds little volume here (0.12 shares per share bought or less,
barely measurable). The thin day's lower cost per share is this market's doing: its fundamentalists pull the price back at a rate that does not
thin with the day, and a slower order leaves them more time. A real thin day usually costs more; what carries over is that the algorithm's clock is
the volume, and on a thin day its own trades are a larger part of that clock.

\section{Implementation-shortfall and close algorithms}

The \omterm{def:mx:benchmark-algorithms:is}{implementation-shortfall algorithm} of \cref{tab:mx:benchmark-algorithms:five} follows chapter~14's trajectory with $\kappa T=4$: it front-loads,
buying 87\% of the order in the first half of the day. Its benchmark is the arrival price, so it is the only algorithm measured against the price
the order actually met; its cost against arrival, 2.9 ticks, has the largest spread of the measured-against-benchmark numbers only because the others'
benchmarks move with the day. The \omterm{def:mx:benchmark-algorithms:is}{target-close algorithm} keeps 30\% for the closing auction (a market-on-close order sent 30 seconds before the
close) and trades the rest over the last nine bins in proportion to the profile. Its tight tracking of the close is bought with the largest impact
cost (6.8 ticks), and its apparent gain against the close is the close it moved.

What the table does not show is the choice between them, which belongs to the portfolio manager: a VWAP benchmark suits an order whose \omterm{def:mx:the-almgren-chriss-framework:urgency}{urgency} is
low and whose trader wants to look like the day's average; implementation shortfall suits an order with alpha that decays; the close suits funds
valued at the close. Chapter~19 measures algorithms against each other with transaction-cost analysis, and chapter~28 builds one that trades each
slice through the placement and routing of chapters 17 and 18.

\section{Tutorial: volume forecasts and benchmark algorithms}\label{tut:mx:benchmark-algorithms:tut}

\begin{tutorial}
\textbf{Goal.} Forecast intraday volume three ways, run the five algorithms in the simulated market, and split VWAP slippage into its schedule and
execution parts.
\textbf{End state:} \cref{fig:mx:benchmark-algorithms:profile,fig:mx:benchmark-algorithms:forecast,fig:mx:benchmark-algorithms:vwap} and the two tables.
\begin{tutsteps}
\item \textbf{Panel.} \texttt{mx\_algos.panel()} and \texttt{flow()}: the activity curves and the volume the simulator prints.
\item \textbf{Forecasts.} \texttt{firm\_algos.Profile}, \texttt{LevelAR.fit(history)}, \texttt{PCAARMA(window).model(j)}; \texttt{forecast\_study()}.
\item \textbf{Agent.} \texttt{Slicer(schedule)} on \texttt{firm.agentmkt}'s \texttt{session}; \texttt{run}, \texttt{day\_stats}, \texttt{fills}.
\item \textbf{VWAP.} \texttt{vwap\_schedules(j, d)}, \texttt{vwap\_study()} with \texttt{firm\_algos.decompose}.
\item \textbf{Participation and the rest.} \texttt{pov\_study()}, \texttt{benchmark\_study()}; draw with \texttt{fig\_algos.py}.
\end{tutsteps}
\textbf{What to change next.} Let the price's volatility rise with volume (news moves prices), which makes the schedule part larger exactly on the
days the forecast is worst; add a second PCA factor; give the \omterm{def:mx:benchmark-algorithms:vwap}{participation algorithm} a cap and a floor on its rate.
\end{tutorial}

\section{Build: benchmark algorithms}\label{bld:mx:benchmark-algorithms:algos}

\begin{build}
\bfield{Purpose} The benchmark schedules and the volume forecasters that chapter~28's \omterm{def:mx:benchmark-algorithms:algo}{execution algorithm}, the transaction-cost analysis of
chapter~19 and the desk of chapter~20 call.
\bfield{Interface} \texttt{profile(history)}, \texttt{Profile(prof, adv)}, \texttt{LevelAR(prof, adv, rho, s\_level, s\_dev)} and
\texttt{LevelAR.fit(history)}, \texttt{PCAARMA(window, r)} and \texttt{.model(j)}, each with \texttt{remaining(observed)};
\texttt{vwap(Q, model, day, dynamic)}, \texttt{twap(Q, bins)}, \texttt{pov(rate, others, own\_counted)}, \texttt{shortfall(Q, bins, kT)} (on
\texttt{firm.acexec}), \texttt{close(Q, model, moc\_share, start)}, \texttt{decompose(q, fills, v, vwaps)}, \texttt{tracking(q, v)}.
\bfield{Rules} Volumes per bin, days $\times$ bins for histories; schedules in shares per bin, summing to the order; a dynamic schedule sees only bins
already finished; slippage per share, positive for a buyer who paid more.
\bfield{Acceptance tests} \texttt{code/firm/algos/tests/}: a static model's dynamic schedule equals its static one; a hot bin tilts the next and a busy
day on profile does not; the level-plus-AR fit recovers its parameters; PCA-ARMA recovers a common shape; the participation steady state; the
decomposition adds up.
\bfield{Stretch} Volume forecasts with a regime for news days; a participation band (chapter~28); an auction-imbalance forecast for the close.
\end{build}

\begin{omsources}
\item H.~Konishi, ``Optimal slice of a VWAP trade'', \emph{Journal of Financial Markets} 5(2), 2002.
\item J.~Białkowski, S.~Darolles and G.~Le~Fol, ``Improving VWAP strategies: a dynamic volume approach'', \emph{Journal of Banking and Finance}
32(9), 2008.
\item A.~Madhavan, ``Volume-weighted average price (VWAP)'', \emph{Encyclopedia of Quantitative Finance}, Wiley, 2010.
\item M.~Humphery-Jenner, ``Optimal VWAP trading under noisy conditions'', \emph{Journal of Banking and Finance} 35(9), 2011.
\item R.~Kissell, \emph{The Science of Algorithmic Trading and Portfolio Management}, Academic Press, 2014.
\end{omsources}

\section{Exercises}

\begin{exercise}[$\star$]\label{exo:mx:benchmark-algorithms:1}
The simulated stock's average profile gives 5.50\% of the day's volume to the first bin and 9.38\% to the last. How many shares of a 15,000-share
static VWAP order go to each?
\end{exercise}

\begin{exercise}[$\star$]\label{exo:mx:benchmark-algorithms:2}
A \omterm{def:mx:benchmark-algorithms:vwap}{participation algorithm} is told to be 10\% of the volume but follows only the other traders' volume, trading 10\% of it. What share of all the
volume does it reach, and what rate on the others' volume would give 10\%?
\end{exercise}

\begin{exercise}[$\star$]\label{exo:mx:benchmark-algorithms:3}
A buyer trades 2 shares in each of two bins, paying 10.10 and 11.00; the market traded 1 and 3 shares there at VWAPs 10.00 and 11.00. Split the
slippage against the day's VWAP into its execution and schedule parts.
\end{exercise}

\begin{exercise}[$\star\star$]\label{exo:mx:benchmark-algorithms:4}
Why does the schedule part explain 99\% of the VWAP slippage's variance on high-volume days but only about 70\% on the other days?
\end{exercise}

\begin{exercise}[$\star\star$]\label{exo:mx:benchmark-algorithms:5}
The \omterm{def:mx:benchmark-algorithms:is}{target-close algorithm} beat the close by 5.9 ticks and had the largest cost against the arrival price. Reconcile the two.
\end{exercise}

\begin{exercise}[$\star\star$]\label{exo:mx:benchmark-algorithms:6}
Three \omterm{def:mx:benchmark-algorithms:vwap}{participation algorithms} at $r=20\%$ buy the same stock, each following the whole volume. At what pace does each trade, relative to the volume
the market would have had without them? What happens at $r=25\%$?
\end{exercise}

\begin{exercise}[$\star\star\star$]\label{exo:mx:benchmark-algorithms:7}
\emph{Coding.} Rerun the forecast study with news on 20\% of stock-days instead of 8\%. How many stock-days are high-volume, and what are the
static curve's and the level-plus-AR schedule's tracking errors on them?
\end{exercise}

\begin{exercise}[$\star\star\star$]\label{exo:mx:benchmark-algorithms:8}
\emph{Find the flaw.} ``Our \omterm{def:mx:benchmark-algorithms:vwap}{VWAP algorithm} finished within 0.6 ticks of the VWAP on average, so it cost our client 0.6 ticks.''
\end{exercise}

\section{Problem: VWAP on a Day That Was Not Average}

\begin{problem}[{Weekend problem --- VWAP on a day that was not average}]\label{pb:mx:benchmark-algorithms:1}
A client asks why its VWAP orders miss the benchmark most on the days the stock is busiest, and whether a dynamic algorithm would help.

\medskip\noindent\textbf{Part I --- The algorithms.}
\begin{enumerate}
\item Define the VWAP, TWAP, participation, implementation-shortfall and \omterm{def:mx:benchmark-algorithms:is}{target-close algorithms} and their benchmarks.
\item What did the \omterm{def:mx:benchmark-algorithms:vwap}{VWAP algorithm}'s buying do to the day's VWAP in the simulated market, and what was its impact cost?
\item Why does the \omterm{def:mx:benchmark-algorithms:is}{target-close algorithm} appear to beat its benchmark?
\item Which measure ranks the five algorithms reliably on eight days, and why?
\end{enumerate}

\medskip\noindent\textbf{Part II --- Forecasting volume.}
\begin{enumerate}[resume]
\item Describe the simulated panel and what makes a high-volume day.
\item State the level-plus-AR model and how it updates during the day.
\item Prove that the schedule part's variance is $s^2\E\sum D_k^2$.
\item Give the tracking errors of the static curve and the level-plus-AR schedule on high-volume and other days.
\item Why does the dynamic PCA-ARMA schedule do worse than the static one on high-volume days?
\end{enumerate}

\medskip\noindent\textbf{Part III --- VWAP in the simulated market.}
\begin{enumerate}[resume]
\item Write the exact split of VWAP slippage into execution and schedule parts.
\item How were the high-volume days chosen, and why do both schedules see the same order flow?
\item Give the standard deviations of the slippage, of its schedule part and of its execution part for both schedules on high-volume days.
\item What is the execution part's mean, and where does it come from?
\item \emph{State the named result}: how much of the VWAP slippage is explained by volume-forecast error, for static and \omterm{def:mx:benchmark-algorithms:dvwap}{dynamic VWAP}, on the
simulator's high-volume days?
\item Is the dynamic schedule's gain reliable? Give the evidence.
\end{enumerate}

\medskip\noindent\textbf{Part IV --- Participation and the answer.}
\begin{enumerate}[resume]
\item Derive the pace $r/(1-r)$ of a \omterm{def:mx:benchmark-algorithms:vwap}{participation algorithm} that counts its own prints.
\item What did a second buyer do to the pace and the cost on the normal day?
\item Why was the pace higher on the thin day?
\item What would you tell the client about \omterm{def:mx:benchmark-algorithms:dvwap}{dynamic VWAP}?
\item In one sentence: what is a benchmark algorithm's slippage a measure of?
\end{enumerate}
\end{problem}

\section{Interview questions}

\begin{interviewq}[$\star$ \iqroles{trader}]\label{iq:mx:benchmark-algorithms:1}
A client wants to buy 5\% of the day's volume in a liquid stock with no view. Which algorithm, and why?
\end{interviewq}

\begin{interviewq}[$\star\star$ \iqroles{researcher}]\label{iq:mx:benchmark-algorithms:2}
How would you forecast the intraday volume curve of a stock, and how would you judge the forecast?
\end{interviewq}

\begin{interviewq}[$\star\star$ \iqroles{researcher}]\label{iq:mx:benchmark-algorithms:3}
Why does a \omterm{def:mx:benchmark-algorithms:vwap}{VWAP algorithm}'s slippage depend on the volume forecast only through cumulative errors?
\end{interviewq}

\begin{interviewq}[$\star\star$ \iqroles{trader}]\label{iq:mx:benchmark-algorithms:4}
Your \omterm{def:mx:benchmark-algorithms:vwap}{participation algorithm} is at 20\% on a stock where another large buyer is also running one. What happens, and what would you do?
\end{interviewq}

\begin{interviewq}[$\star\star$ \iqroles{developer}]\label{iq:mx:benchmark-algorithms:5}
Design the interface between a schedule and the child-order layer of an \omterm{def:mx:benchmark-algorithms:algo}{execution algorithm}.
\end{interviewq}

\begin{interviewq}[$\star\star\star$ \iqroles{bank}]\label{iq:mx:benchmark-algorithms:6}
A client complains that your \omterm{def:mx:benchmark-algorithms:is}{target-close algorithm} beat the close but the fund's performance suffered. What could have happened?
\end{interviewq}
